\section{AHSME 1952}
        \begin{enumerate}
        	\item (\textit{AHSME 1952}) If the radius of a circle is a rational number, its area is given by a number which is:


 $\textbf{(A)\ } \text{rational}  \qquad \textbf{(B)\ } \text{irrational} \qquad \textbf{(C)\ } \text{integral} \qquad \textbf{(D)\ } \text{a perfect square }\qquad \textbf{(E)\ } \text{none of these}$

 \textbf{Answer: }B



\textbf{Solution}

Let the radius of the circle be the common fraction $\frac{a}{b}.$ Then the area of the circle is $\pi \cdot \frac{a^2}{b^2}.$Because $\pi$ is irrational and $\frac{a^2}{b^2}$ is rational, their product must be irrational. The answer is $\boxed{B}.$




\textbf{Solution 2}

The phrasing of the problem makes it clear that the rule would apply to all rational radii. So let r be equal to one. $\pi$ times 1 squared is equal to $\pi$ which is irrational. Therefore, the answer is $\boxed{B}.$



	\item (\textit{AHSME 1952}) Two high school classes took the same test. One class of $20$ students made an average grade of $80\%$; the other class of $30$ students made an average grade of $70\%$. The average grade for all students in both classes is:


 $\textbf{(A)}\ 75\%\qquad \textbf{(B)}\ 74\%\qquad \textbf{(C)}\ 72\%\qquad \textbf{(D)}\ 77\%\qquad \textbf{(E)\ }\text{none of these}$

 \textbf{Answer: }B



\textbf{Solution}

The desired average can be found by dividing the total number of points earned by the total number of students. There are $20\cdot 80+30\cdot 70=3700$ points earned and $20+30=50$ students. Thus, our answer is $\frac{3700}{50}$, or $\boxed{\textbf{(B)}\ 74\%}$.

 


	\item (\textit{AHSME 1952}) The expression $a^3-a^{-3}$ equals: 


 $\textbf{(A) \ }\left(a-\frac{1}{a}\right)\left(a^2+1+\frac{1}{a^2}\right) \qquad \textbf{(B) \ }\left(\frac{1}{a}-a\right)\left(a^2-1+\frac{1}{a^2}\right) \qquad \textbf{(C) \ }\left(a-\frac{1}{a}\right)\left(a^2-2+\frac{1}{a^2}\right) \qquad$
$\textbf{(D) \ }\left(\frac{1}{a}-a\right)\left(\frac{1}{a^2}+1+a^2\right) \qquad \textbf{(E) \ }\text{none of these}$

 \textbf{Answer: }A



\textbf{Solution}

Recall that $x^3-y^3=(x-y)(x^2+xy+y^2)$. Letting $a=x$ and $a^{-1}=y$, we find that $a^3-a^{-3}= \boxed{\textbf{(A)}\ \left(a-\frac{1}{a}\right)\left(a^2+1+\frac{1}{a^2}\right)}$.

 


	\item (\textit{AHSME 1952}) The cost $C$ of sending a parcel post package weighing $P$ pounds, $P$ an integer, is $10$ cents for the first pound and $3$ cents for each additional pound. The formula for the cost is:


 $\textbf{(A) \ }C=10+3P \qquad \textbf{(B) \ }C=10P+3 \qquad \textbf{(C) \ }C=10+3(P-1) \qquad$
$\textbf{(D) \ }C=9+3P \qquad \textbf{(E) \ }C=10P-7$

 \textbf{Answer: }C



\textbf{Solution}

We know that the first pound has a constant price, $10$ cents, and that the other $P-1$ pounds cost $3$ cents apiece. This leaves us with $\boxed{\textbf{(C)}\ C=10+3(P-1)}$.

 


	\item (\textit{AHSME 1952}) The points $(6,12)$ and $(0,-6)$ are connected by a straight line. Another point on this line is:


 $\textbf{(A) \ }(3,3)  \qquad \textbf{(B) \ }(2,1) \qquad \textbf{(C) \ }(7,16) \qquad \textbf{(D) \ }(-1,-4) \qquad \textbf{(E) \ }(-3,-8)$

 \textbf{Answer: }A



\textbf{Solution}

The slope of this line is $\frac{y_2-y_1}{x_2-x_1}=\frac{12+6}{6-0}=3$. Hence, its equation is $y=3x-6$. The only given point which satisfies these conditions is $\boxed{\textbf{(A)}\ (3,3)}$.

 


	\item (\textit{AHSME 1952}) The difference of the roots of $x^2-7x-9=0$ is:


 $\textbf{(A) \ }+7  \qquad \textbf{(B) \ }+\frac{7}{2} \qquad \textbf{(C) \ }+9 \qquad \textbf{(D) \ }2\sqrt{85} \qquad \textbf{(E) \ }\sqrt{85}$

 \textbf{Answer: }E



\textbf{Solution}

Denote the $2$ roots of this quadratic as $r_1$ and $r_2$. Note that $(r_1-r_2)^2=(r_1+r_2)^2-4r_1r_2$. By \href{https://artofproblemsolving.com/wiki/index.php?title=Vieta\%27s}{Vieta's Formula's}, $r_1+r_2=7$, and $r_1r_2=-9$. Thus, $r_1-r_2=\sqrt{49+4\cdot 9}=\boxed{\textbf{(E)}\ \sqrt{85}}$.

 


	\item (\textit{AHSME 1952}) When simplified, $(x^{-1}+y^{-1})^{-1}$ is equal to:


 $\textbf{(A) \ }x+y  \qquad \textbf{(B) \ }\frac{xy}{x+y} \qquad \textbf{(C) \ }xy \qquad \textbf{(D) \ }\frac{1}{xy} \qquad \textbf{(E) \ }\frac{x+y}{xy}$

 \textbf{Answer: }B



\textbf{Solution}

$(x^{-1}+y^{-1})^{-1}=\left(\frac{1}{x}+\frac{1}{y}\right)^{-1}=\left(\frac{x+y}{xy}\right)^{-1}=\boxed{\textbf{(B)}\ \frac{xy}{x+y}}$.

 


	\item (\textit{AHSME 1952}) Two equal circles in the same plane cannot have the following number of common tangents.


 $\textbf{(A) \ }1  \qquad \textbf{(B) \ }2 \qquad \textbf{(C) \ }3 \qquad \textbf{(D) \ }4 \qquad \textbf{(E) \ }\text{none of these}$

 \textbf{Answer: }A



\textbf{Solution}

Two congruent coplanar circles will either be tangent to one another (resulting in $3$ common tangents), intersect one another (resulting in $2$ common tangents), or be separate from one another (resulting in $4$ common tangents).Having only $\boxed{\textbf{(A)}\ 1}$ common tangent is impossible, unless the circles are non-congruent and internally tangent.

 
                    \begin{figure}[H]
                        \centering
                        \includegraphics[width=0.5\textwidth]{Solution/AHSME_I_1952_8_Solution_Figure_1.png}
                    \end{figure}
                    



	\item (\textit{AHSME 1952}) If $m=\frac{cab}{a-b}$, then $b$ equals:


 $\textbf{(A) \ }\frac{m(a-b)}{ca}  \qquad \textbf{(B) \ }\frac{cab-ma}{-m} \qquad \textbf{(C) \ }\frac{1}{1+c} \qquad \textbf{(D) \ }\frac{ma}{m+ca} \qquad \textbf{(E) \ }\frac{m+ca}{ma}$

 \textbf{Answer: }D



\textbf{Solution}

$m=\frac{cab}{a-b}\implies ma-mb=cab\implies ma=b(m+ac)\implies b=\boxed{\textbf{(D)}\ \frac{ma}{m+ca}}$



	\item (\textit{AHSME 1952}) An automobile went up a hill at a speed of $10$ miles an hour and down the same distance at a speed of $20$ miles an hour. The average speed for the round trip was:


 $\textbf{(A) \ }12\frac{1}{2}\text{mph}  \qquad \textbf{(B) \ }13\frac{1}{3}\text{mph} \qquad \textbf{(C) \ }14\frac{1}{2}\text{mph} \qquad \textbf{(D) \ }15\text{mph} \qquad \textbf{(E) \ }\text{none of these}$

 \textbf{Answer: }B



\textbf{Solution}

Let $r$, $t$, and $d$ represent the rate, time, and distance (respectively) of each trip. We know that $rt=d$, or $r_1t_1=r_2t_2$. Because it is given that $2r_1=r_2$, $t_1=2t_2$. Furthermore, $d=10t_1=20t_2$. The average speed can be expressed as $\frac{2d}{t_1+t_2}=\frac{2(20t_2)}{2t_2+t_2}=\frac{40}{3}=\boxed{\textbf{(B)}\ 13\frac{1}{3}\text{ mph}}$.

 


	\item (\textit{AHSME 1952}) If $y=f(x)=\frac{x+2}{x-1}$, then it is incorrect to say:


 $\textbf{(A)\ }x=\frac{y+2}{y-1}\qquad\textbf{(B)\ }f(0)=-2\qquad\textbf{(C)\ }f(1)=0\qquad$
$\textbf{(D)\ }f(-2)=0\qquad\textbf{(E)\ }f(y)=x$

 \textbf{Answer: }C



\textbf{Solution}

$f(1)=\frac{3}{0}$, which is undefined. Hence, the (in)correct answer is $\boxed{\textbf{(C)}\ f(1)=0}$.We can verify that the other statements are valid.

 $\textbf{(A).\ } y=\frac{x+2}{x-1}\implies xy-y=x+2\implies x(y-1)=y+2\implies x=\frac{y+2}{y-1}$
$\textbf{(B).\ } f(0)=\frac{(0)+2}{(0)-1}=-2$
$\textbf{(D).\ } f(-2)=\frac{(-2)+2}{(-2)-1}=\frac{0}{-3}=0$
$\textbf{(E).\ } f(y)=\frac{\frac{x+2}{x-1}+2}{\frac{x+2}{x-1}-1}=\frac{3x}{3}=x$



	\item (\textit{AHSME 1952}) The sum to infinity of the terms of an infinite geometric progression is $6$. The sum of the first two terms is $4\frac{1}{2}$. The first term of the progression is:


 $\textbf{(A) \ }3 \text{ or } 1\frac{1}{2}  \qquad \textbf{(B) \ }1 \qquad \textbf{(C) \ }2\frac{1}{2} \qquad \textbf{(D) \ }6 \qquad \textbf{(E) \ }9\text{ or }3$

 \textbf{Answer: }E



\textbf{Solution}

This geometric sequence can be written as $a+ar+ar^2+ar^3+\cdots$. We are given that $a+ar=4\frac{1}{2}$. Using the \href{https://artofproblemsolving.com/wiki/index.php?title=Geometric\_sequence\#Infinite\_Geometric\_Sequences}{formula} for the sum of an infinite geometric series, we know that $\frac{a}{1-r}=6$. Solving for $r$ in the second equation, we find that $r=\frac{6-a}{6}$. Plugging this into the first equation results in $a^2-12a+27=0$, which can be factored as $(a-3)(a-9)=0$. Hence, $a$ equals $\boxed{\textbf{(E)}\ 9 \text{ or }3}$.

 


	\item (\textit{AHSME 1952}) The function $x^2+px+q$ with $p$ and $q$ greater than zero has its minimum value when:


 $\textbf{(A) \ }x=-p  \qquad \textbf{(B) \ }x=\frac{p}{2} \qquad \textbf{(C) \ }x=-2p \qquad \textbf{(D) \ }x=\frac{p^2}{4q} \qquad$
$\textbf{(E) \ }x=\frac{-p}{2}$

 \textbf{Answer: }E



\textbf{Solution}

The minimum value of this parabola is found at its turning point, on the line $\boxed{\textbf{(E)}\ x=\frac{-p}{2}}$.Indeed, the turning point of any function of the form $ax^2+bx+c$ has an x-coordinate of $\frac{-b}{2a}$. This can be seen at the average of the quadratic's two roots (whose sum is $\frac{-b}{a}$) or (using calculus) as the value of its derivative set equal to $0$.

 


	\item (\textit{AHSME 1952}) A house and store were sold for $\textdollar 12,000$ each. The house was sold at a loss of $20\%$ of the cost, and the store at a gain of $20\%$ of the cost. The entire transaction resulted in:


 $\textbf{(A) \ }\text{no loss or gain}  \qquad \textbf{(B) \ }\text{loss of }\textdollar 1000 \qquad \textbf{(C) \ }\text{gain of }\textdollar 1000 \qquad \textbf{(D) \ }\text{gain of }\textdollar 2000 \qquad \textbf{(E) \ }\text{none of these}$

 \textbf{Answer: }B



\textbf{Solution}

Denote the original price of the house and the store as $h$ and $s$, respectively. It is given that $\frac{4h}{5}=\textdollar 12,000$, and that $\frac{6s}{5}=\textdollar 12,000$. Thus, $h=\textdollar 15,000$, $s=\textdollar10,000$, and $h+s=\textdollar25,000$. This value is $\textdollar1000$ higher than the current price of the property, $2\cdot \textdollar12,000$. Hence, the transaction resulted in a $\boxed{\textbf{(B)}\ \text{loss of }\textdollar1000}$.

 


	\item (\textit{AHSME 1952}) The sides of a triangle are in the ratio $6:8:9$. Then:


 $\textbf{(A) \ }\text{the triangle is obtuse}$
$\textbf{(B) \ }\text{the angles are in the ratio }6:8:9$
$\textbf{(C) \ }\text{the triangle is acute}$
$\textbf{(D) \ }\text{the angle opposite the largest side is double the angle opposite the smallest side}$
$\textbf{(E) \ }\text{none of these}$

 \textbf{Answer: }C



\textbf{Solution}

In a triangle with sides $a$, $b$, and $c$, where $a\le b\le c$, $\bigtriangleup ABC$ is acute if $a^2+b^2>c^2$, right if $a^2+b^2=c^2$, and obtuse if $a^2+b^2<c^2$. If $c$ were equal to $10$, $\bigtriangleup ABC$ would be right, as a multiple of the Pythagorean triple $(3,4,5)$. Because $c<10$, $\boxed{\textbf{(C)}\ \text{the triangle is acute}}$.

 


	\item (\textit{AHSME 1952}) If the base of a rectangle is increased by $10\%$ and the area is unchanged, then the altitude is decreased by:


 $\textbf{(A) \ }9\%  \qquad \textbf{(B) \ }10\% \qquad \textbf{(C) \ }11\% \qquad \textbf{(D) \ }11\frac{1}{9}\% \qquad \textbf{(E) \ }9\frac{1}{11}\%$

 \textbf{Answer: }E



\textbf{Solution}

$b\cdot a=\frac{11}{10}b\cdot xa\implies x=\frac{10}{11}$. Hence, the altitude is decreased by $1-x=\frac{1}{11}=\boxed{\textbf{(E)}\ 9\frac{1}{11}\%}$.

 


	\item (\textit{AHSME 1952}) A merchant bought some goods at a discount of $20\%$ of the list price. He wants to mark them at such a price that he can give a discount of $20\%$ of the marked price and still make a profit of $20\%$ of the selling price. The percent of the list price at which he should mark them is:


 $\textbf{(A) \ }20  \qquad \textbf{(B) \ }100 \qquad \textbf{(C) \ }125 \qquad \textbf{(D) \ }80 \qquad \textbf{(E) \ }120$

 \textbf{Answer: }C



\textbf{Solution}

Let $C$ represent the cost of the goods, and let $L$, $S$, and $M$ represent the list, selling, and marked prices of the goods, respectively. Hence, we have three equations, which we need to manipulate in order to relate $M$ and $L$:

 $C=\frac{4}{5}L$
$S=C+\frac{1}{5}S$
$S=\frac{4}{5}M$

 We find that $M=\frac{5}{4}L$. Hence, the percent of the list price which should be marked is $\boxed{\textbf{(C)}\ 125}$.

 


	\item (\textit{AHSME 1952}) $\log p+\log q=\log(p+q)$ only if:


 $\textbf{(A) \ }p=q=\text{zero}  \qquad \textbf{(B) \ }p=\frac{q^2}{1-q} \qquad \textbf{(C) \ }p=q=1 \qquad$
$\textbf{(D) \ }p=\frac{q}{q-1} \qquad \textbf{(E) \ }p=\frac{q}{q+1}$

 \textbf{Answer: }D



\textbf{Solution}

$\log p+\log q=\log(p+q)\implies \log pq=\log(p+q)\implies pq=p+q\implies \boxed{\textbf{(D)}\ p=\frac{q}{q-1}}$



	\item (\textit{AHSME 1952}) Angle $B$ of triangle $ABC$ is trisected by $BD$ and $BE$ which meet $AC$ at $D$ and $E$ respectively. Then:


 $\textbf{(A) \ }\frac{AD}{EC}=\frac{AE}{DC}  \qquad \textbf{(B) \ }\frac{AD}{EC}=\frac{AB}{BC} \qquad \textbf{(C) \ }\frac{AD}{EC}=\frac{BD}{BE} \qquad$
$\textbf{(D) \ }\frac{AD}{EC}=\frac{(AB)(BD)}{(BE)(BC)} \qquad \textbf{(E) \ }\frac{AD}{EC}=\frac{(AE)(BD)}{(DC)(BE)}$

 \textbf{Answer: }D



\textbf{Solution}


                    \begin{figure}[H]
                        \centering
                        \includegraphics[width=0.5\textwidth]{Solution/AHSME_I_1952_19_Solution_Figure_1.png}
                    \end{figure}
                    Using the \href{https://artofproblemsolving.com/wiki/index.php?title=Angle\_Bisector\_Theorem}{Angle Bisector Theorem} on $\angle ABE$ and $\angle DBC$, we find that $\frac{AD}{DE}=\frac{AB}{BE}$ and $\frac{DE}{EC}=\frac{BD}{BC}$. Rewriting this to better fit our answer choices gives $AD=\frac{(DE)(AB)}{BE}$ and $EC=\frac{(DE)(BC)}{BD}$. Hence, $\frac{AD}{EC}=\frac{(DE)(AB)}{BE}\cdot \frac{BD}{(DE)(BC)}\implies \boxed{\textbf{(D)}\ \frac{AD}{EC}=\frac{(AB)(BD)}{(BE)(BC)}}$.

 


	\item (\textit{AHSME 1952}) If $\frac{x}{y}=\frac{3}{4}$, then the incorrect expression in the following is:


 $\textbf{(A) \ }\frac{x+y}{y}=\frac{7}{4}  \qquad \textbf{(B) \ }\frac{y}{y-x}=\frac{4}{1} \qquad \textbf{(C) \ }\frac{x+2y}{x}=\frac{11}{3} \qquad$
$\textbf{(D) \ }\frac{x}{2y}=\frac{3}{8} \qquad \textbf{(E) \ }\frac{x-y}{y}=\frac{1}{4}$

 \textbf{Answer: }E



\textbf{Solution}

Let's consider each answer choice.

 $\boxed{A.} \quad \frac{x+y}{y}=\frac{x}{y}+\frac{y}{y}=\frac{3}{4}+1=\boxed{\frac{7}{4}}$
$\boxed{B.} \quad \frac{y-x}{y}=\frac{y}{y}-\frac{x}{y}=1-\frac{3}{4}=\frac{1}{4}\implies \left(\frac{y-x}{y}\right)^{-1}=\frac{y}{y-x}=\boxed{\frac{4}{1}}$
$\boxed{C.} \quad \frac{x+2y}{x}=\frac{x}{x}+\frac{2y}{x}=1+2\cdot \frac{4}{3}=\boxed{\frac{11}{3}}$
$\boxed{D.} \quad \frac{x}{2y}=\frac{1}{2}\cdot \frac{3}{4}=\boxed{\frac{3}{8}}$
$\boxed{E.} \quad \frac{x-y}{y}=\frac{x}{y}-\frac{y}{y}=\frac{3}{4}-1=\boxed{-\frac{1}{4}}\neq \frac{1}{4}$

 Hence, the incorrect expression is $\boxed{\textbf{(E)}\ \frac{x-y}{y}=\frac{1}{4}}$.

 Note: If you're bad at fraction manipulation, you can plug in $x=3$ and $y=4$ and find which expression is false.

 


	\item (\textit{AHSME 1952}) The sides of a regular polygon of $n$ sides, $n>4$, are extended to form a star. The number of degrees at each point of the star is:


 $\textbf{(A) \ }\frac{360}{n}  \qquad \textbf{(B) \ }\frac{(n-4)180}{n} \qquad \textbf{(C) \ }\frac{(n-2)180}{n} \qquad$
$\textbf{(D) \ }180-\frac{90}{n} \qquad \textbf{(E) \ }\frac{180}{n}$

 \textbf{Answer: }B



\textbf{Solution}


                    \begin{figure}[H]
                        \centering
                        \includegraphics[width=0.5\textwidth]{Solution/AHSME_I_1952_21_Solution_Figure_1.png}
                    \end{figure}
                    

 The measure of each angle, in degrees, of an $n$-sided regular polygon is $\frac{(n-2)180}{n}$. Hence, the two base angles of each triangle formed measure $180-\frac{(n-2)180}{n}$ degrees, and the vertex angle measures $180-2\left(180-\frac{(n-2)180}{n}\right)$ degrees. This simplifies to $\frac{180n-720}{n}$, or $\boxed{\textbf{(B)}\ \frac{(n-4)180}{n}}$.

 


	\item (\textit{AHSME 1952}) On hypotenuse $AB$ of a right triangle $ABC$ a second right triangle $ABD$ is constructed with hypotenuse $AB$. If $\overline{BC}=1$, $\overline{AC}=b$, and $\overline{AD}=2$, then $\overline{BD}$ equals:


 $\textbf{(A) \ }\sqrt{b^2+1}  \qquad \textbf{(B) \ }\sqrt{b^2-3} \qquad \textbf{(C) \ }\sqrt{b^2+1}+2 \qquad$
$\textbf{(D) \ }b^2+5 \qquad \textbf{(E) \ }\sqrt{b^2+3}$

 \textbf{Answer: }B



\textbf{Solution}


                    \begin{figure}[H]
                        \centering
                        \includegraphics[width=0.5\textwidth]{Solution/AHSME_I_1952_22_Solution_Figure_1.png}
                    \end{figure}
                    Let $\overline{BD}=x$. According to the \href{https://artofproblemsolving.com/wiki/index.php?title=Pythagorean\_Theorem}{Pythagorean Theorem}, $x^2+4=b^2+1$. Hence, $x=\boxed{\textbf{(B)}\ \sqrt{b^2-3}}$.

 


	\item (\textit{AHSME 1952}) If $\frac{x^2-bx}{ax-c}=\frac{m-1}{m+1}$ has roots which are numerically equal but of opposite signs, the value of $m$ must be:


 $\textbf{(A) \ }\frac{a-b}{a+b} \qquad \textbf{(B) \ }\frac{a+b}{a-b} \qquad \textbf{(C) \ }c \qquad \textbf{(D) \ }\frac{1}{c} \qquad \textbf{(E) \ }1$

 \textbf{Answer: }A



\textbf{Solution}

Cross-multiplying, we find that $(m+1)x^2-(bm+am+b-a)x+c(m-1)=0$. Because the roots of this quadratic are additive inverses, their sum is $0$. According to \href{https://artofproblemsolving.com/wiki/index.php?title=Vieta\%27s\_Formulas}{Vieta's Formulas}, $\frac{bm+am+b-a}{m+1}=0$, or $m(a+b)=a-b$. Hence, $m=\boxed{\textbf{(A)}\ \frac{a-b}{a+b}}$.

 


	\item (\textit{AHSME 1952}) In the figure, it is given that angle $C = 90^{\circ}$,$\overline{AD} = \overline{DB}$,$DE \perp AB$, $\overline{AB} = 20$, and $\overline{AC} = 12$. The area of quadrilateral $ADEC$ is: 

                    \begin{figure}[H]
                        \centering
                        \includegraphics[width=0.5\textwidth]{Problem/AHSME_I_1952_24_Problem_Figure_1.png}
                    \end{figure}
                    
$\textbf{(A)}\ 75\qquad\textbf{(B)}\ 58\frac{1}{2}\qquad\textbf{(C)}\ 48\qquad\textbf{(D)}\ 37\frac{1}{2}\qquad\textbf{(E)}\ \text{none of these}$

 \textbf{Answer: }B



\textbf{Solution}

$[ADEC]=[BCA]-[BDE]$ 

 Note that, as right triangles sharing $\angle B$, $\triangle BDE \sim \triangle BCA$. 

 $\overline{BD}=\frac{20}{2}=10$

 Because the sides of $\triangle BCA$ are in the ratio $3:4:5$, $\overline{BC}=16$.

 The sides of our triangles are in the ratio $\frac{10}{16}=\frac{5}{8}$, and the ratio of their areas is $\left(\frac{5}{8}\right)^2=\frac{25}{64}$.

 $[BCA]=\frac{12 \cdot 16}{2}=96$, and $[BDE]=\frac{96 \cdot 25}{64}=\frac{75}{2}$.

 $[ADEC]=\frac{192-75}{2}=\boxed{\textbf{(B)}\ 58\frac{1}{2}}$



	\item (\textit{AHSME 1952}) A powderman set a fuse for a blast to take place in $30$ seconds. He ran away at a rate of $8$ yards per second. Sound travels at the rate of $1080$ feet per second. When the powderman heard the blast, he had run approximately:


 $\textbf{(A)}\ \text{200 yd.}\qquad\textbf{(B)}\ \text{352 yd.}\qquad\textbf{(C)}\ \text{300 yd.}\qquad\textbf{(D)}\ \text{245 yd.}\qquad\textbf{(E)}\ \text{512 yd.}$

 \textbf{Answer: }D



\textbf{Solution}

Let $p(t)=24t$ be the number of feet the powderman is from the blast at $t$ seconds after the fuse is lit, and let $q(t)=1080t-32400$ be the number of feet the sound has traveled. We want to solve for $p(t)=q(t)$.
 \[24t=1080t-32400\]
 \[1056t=32400\]
 \[t=\frac{32400}{1056}\]
 \[t=\frac{675}{22}=30.6\overline{81}\]The number of \textbf{yards} the powderman is from the blast at time $t$ is $\frac{24t}3=8t$, so the answer is $8(30.6\overline{81})$, which is about $245$ yards. $\boxed{\textbf{(D)}}$



	\item (\textit{AHSME 1952}) If $\left(r+\frac1r\right)^2=3$, then $r^3+\frac1{r^3}$ equals


 $\textbf{(A)}\ 1\qquad\textbf{(B)}\ 2\qquad\textbf{(C)}\ 0\qquad\textbf{(D)}\ 3\qquad\textbf{(E)}\ 6$

 \textbf{Answer: }C



\textbf{Solution}

We know $r+\frac1r=\sqrt3$. Cubing this gives $r^3+3r+\frac3r+\frac1{r^3}=3\sqrt3$. But $3r+\frac3r=3\left(r+\frac1r\right)=3\sqrt3$, so subtracting this from the first equation gives$r^3+\frac1{r^3}=\boxed{0\textbf{ (C)}}$. (Actually, $r+\frac1r$ could have been equal to $-\sqrt3$ instead of $\sqrt3$, but this would have led to the same answer. Also, this answer implies that $r^6=-1$, which means that $r$ is a complex number.)

 


	\item (\textit{AHSME 1952}) The ratio of the perimeter of an equilateral triangle having an altitude equal to the radius of a circle, to the perimeter of an equilateral triangle inscribed in the circle is:


 $\textbf{(A)}\ 1:2\qquad\textbf{(B)}\ 1:3\qquad\textbf{(C)}\ 1:\sqrt3\qquad\textbf{(D)}\ \sqrt3:2 \qquad\textbf{(E)}\ 2:3$

 \textbf{Answer: }E



\textbf{Solution}

If the radius of the circle is $r$, then the perimeter of the first triangle is $3\left(\frac{2r}{\sqrt3}\right)=2r\sqrt3$, and the perimeter of the second is $3r\sqrt3$. So the ratio is $\boxed{\frac23{\textbf{ (E)}}}$.

 


	\item (\textit{AHSME 1952}) In the table shown, the formula relating $x$ and $y$ is:


 $\begin{tabular}{|l|l|l|l|l|l|} \hline x & 1 & 2 & 3 & 4 & 5 \\ \hline y & 3 & 7 & 13 & 21 &  31 \\ \hline \end{tabular}$
$\textbf{(A)}\ y=4x-1  \qquad \textbf{(B)}\ y=x^{3}-x^{2}+x+2 \qquad \textbf{(C)}\ y=x^2+x+1 \\ \qquad \textbf{(D)}\ y=(x^2+x+1)(x-1) \qquad \textbf{(E)}\ \text{None of these}$

 \textbf{Answer: }C



\textbf{Solution}

A simple method of solving the problem is trying each of the answer choices. One can plug in values of x and y for each, because many x-values and their single corresponding y-values are given.

 
            \begin{enumerate}
        \item Choice A works for (1,3) and (2,7) but fails to work on (3,13) because 3*4=12 and 12-1=11, not 13.
        \item Choice B doesn't work for (2,7) because it would be 8-4+2+2=8, not 7.
        \item Choice C actually works for all five pairs, being 1+1+1=3, 4+2+1=7, 9+3+1=13, 16+4+1=21, and 25+5+1=31.
            \end{enumerate}
            Thus, the answer is $\boxed{\textbf{(C)}}$.

 


	\item (\textit{AHSME 1952}) In a circle of radius $5$ units, $CD$ and $AB$ are perpendicular diameters. A chord $CH$ cutting $AB$ at $K$ is $8$ units long.  The diameter $AB$ is divided into two segments whose dimensions are:


 $\textbf{(A)}\ 1.25, 8.75  \qquad \textbf{(B)}\ 2.75, 7.25\qquad \textbf{(C)}\ 2,8 \qquad \textbf{(D)} \ 4,6 \qquad \textbf{(E)} \text{None of these}$

 \textbf{Answer: }A



\textbf{Solution}

Let the intersection of $CH$ and $AB$ be $N$, $O$ be the center of the circle, $ON = a$ and $CN = x$. By power of a point on $N$, we have
 \[(5+a)(5-a) = 25-a^2 = x(8-x).\]$\triangle CON$ is a right triangle, so we also know that $x^2 = a^2 + 25$, thus
 \[25 - a^2 = 25-(x^2-25) = 50-x^2 = 8x-x^2 \Rightarrow x = \frac{50}{8} = \frac{25}{4}.\]It follows that $a = \sqrt{x^2 - 25} = 5\sqrt{\frac{25}{16}-1} = \frac{15}{4}$.

 $BN = 5 + \frac{15}{14} = 8.75$, so the answer is $A$.

 


	\item (\textit{AHSME 1952}) When the sum of the first ten terms of an arithmetic progression is four times the sum of the first five terms, the ratio of the first term to the common difference is: 


 $\textbf{(A)}\ 1: 2 \qquad \textbf{(B)}\ 2: 1 \qquad \textbf{(C)}\ 1: 4 \qquad \textbf{(D)}\ 4: 1 \qquad \textbf{(E)}\ 1: 1$

 \textbf{Answer: }A



\textbf{Solution}

Let our first term be $a$ and our common difference be $d$. Thus, the first few terms of the sequence are $a$, $a + d$, $a + 2d$, ...

 The sum of the first 5 terms is
 \[a + (a + d) + (a + 2d) + ...  + (a + 4d) = 5a + 10d\]The sum of the first 10 terms is 
 \[a + (a + d) + (a + 2d) + ... + (a + 9d) = 10a + 45d\]We are told that the latter is four times the former. Hence, 
 \[4(5a + 10d) = 10a + 45d\]
 \[20a + 40d = 10a + 45d\]
 \[10a = 5d\]
 \[\frac{a}{d}  = \frac{1}{2}\]

 Therefore, our answer is $\fbox{(A) 1:2}$



	\item (\textit{AHSME 1952}) Given $12$ points in a plane no three of which are collinear, the number of lines they determine is: 


 $\textbf{(A)}\ 24 \qquad \textbf{(B)}\ 54 \qquad \textbf{(C)}\ 120 \qquad \textbf{(D)}\ 66 \qquad \textbf{(E)}\ \text{none of these}$

 \textbf{Answer: }D



\textbf{Solution}

Since no three points are collinear, every two points must determine a distinct line. Thus, there are  $\dbinom{12}{2} = \frac{12\cdot11}{2} = 66$ lines. 

 Therefore, the answer is $\fbox{(D) 66}$



	\item (\textit{AHSME 1952}) $K$ takes $30$ minutes less time than $M$ to travel a distance of $30$ miles. $K$ travels $\frac {1}{3}$ mile per hour faster than $M$. If $x$ is $K$'s rate of speed in miles per hours, then $K$'s time for the distance is: 


 $\textbf{(A)}\ \dfrac{x + \frac {1}{3}}{30} \qquad \textbf{(B)}\ \dfrac{x - \frac {1}{3}}{30} \qquad \textbf{(C)}\ \frac{30}{x+\frac{1}{3}}\qquad \textbf{(D)}\ \frac{30}{x}\qquad \textbf{(E)}\ \frac{x}{30}$

 \textbf{Answer: }D



\textbf{Solution}

Using the formula d=rt, and setting d=30 and r=x, we can easily see that the answer is $\fbox{D}$.Note that the first sentence is irrelevant.

 


	\item (\textit{AHSME 1952}) A circle and a square have the same perimeter. Then: 


 $\text{(A) their areas are equal}\qquad\\ \text{(B) the area of the circle is the greater} \qquad\\ \text{(C) the area of the square is the greater} \qquad\\ \text{(D) the area of the circle is } \pi \text{ times the area of the square}\qquad\\ \text{(E) none of these}$

 \textbf{Answer: }B



\textbf{Solution}

Assume that the side length of the square is 4 units. Then the perimeter for both shapes is is 4 * 4 = 16 units. Since the perimeter of a circle is 2$\pi$r, then 2$\pi$r =16 and r = 8/$\pi$, which is about 2.5. The area of a circle is $\pi$r\textasciicircum{}2 so (2.5\textasciicircum{}2)*$\pi$ = 6.25$\pi$ or about 19.6. This is more than the area of the square, which is 4\textasciicircum{}2 =16, so the answer is $\fbox{B}$.

 


	\item (\textit{AHSME 1952}) The price of an article was increased $p\%$. Later the new price was decreased $p\%$. If the last price was one dollar, the original price was: 


 $\textbf{(A)}\ \frac{1-p^2}{200}\qquad \textbf{(B)}\ \frac{\sqrt{1-p^2}}{100}\qquad \textbf{(C)}\ \text{one dollar}\qquad\\ \textbf{(D)}\ 1-\frac{p^2}{10000-p^2}\qquad \textbf{(E)}\ \frac{10000}{10000-p^2}$

 \textbf{Answer: }E



\textbf{Solution}

Converting the word problem in to an equation, we take x as the initial amount, we get 
 \[x(1+\dfrac{p}{100})(1-\dfrac{p}{100}) = 1\] Using, $(a+b)(a-b) = a^{2} - b^{2}$. We now simplify it to, 
 \[x(1-\dfrac{p^{2}}{10000})=1\] Simplifying to, 
 \[x(\dfrac{10000-p^{2}}{10000}) = 1\] Dividing both sides by $\dfrac{10000-p^{2}}{10000}$, we get the value of x. $\boxed{\textbf{(E)}\dfrac{10000}{10000-p^{2}}}$



	\item (\textit{AHSME 1952}) With a rational denominator, the expression $\frac {\sqrt {2}}{\sqrt {2} + \sqrt {3} - \sqrt {5}}$ is equivalent to: 


 $\textbf{(A)}\ \frac {3 + \sqrt {6} + \sqrt {15}}{6} \qquad \textbf{(B)}\ \frac {\sqrt {6} - 2 + \sqrt {10}}{6} \qquad \textbf{(C)}\ \frac{2+\sqrt{6}+\sqrt{10}}{10} \qquad\\ \textbf{(D)}\ \frac {2 + \sqrt {6} - \sqrt {10}}{6} \qquad \textbf{(E)}\ \text{none of these}$

 \textbf{Answer: }A



\textbf{Solution}

First let $k=\sqrt2+\sqrt3$. Then our original expression becomes $\frac{\sqrt2}{k-\sqrt5}$. Rationalizing the denominator, we have: \begin{align*}\frac{\sqrt2}{k-\sqrt5}&=\frac{\sqrt2(k+\sqrt5)}{(k-\sqrt5)(k+\sqrt5)}=\frac{\sqrt{2}k+\sqrt{10}}{k^2-(\sqrt5)^2}=\frac{\sqrt{2}k+\sqrt{10}}{k^2-5} \\&=\frac{\sqrt{2}(\sqrt{2}+\sqrt{3})+\sqrt{10}}{(\sqrt{2}+\sqrt{3})^2-5}= \frac{2+\sqrt{6}+\sqrt{10}}{2\sqrt{6}+2+3-5}=\frac{2+\sqrt6+\sqrt{10}}{2\sqrt{6}} \\&=\frac{2\sqrt{6}+6+\sqrt{60}}{12}=\boxed{\textbf{(A) }\frac{3+\sqrt{6}+\sqrt{15}}{6}}.\end{align*}


 (reworked by Technodoggo)

 


	\item (\textit{AHSME 1952}) To be continuous at $x = - 1$, the value of $\frac {x^3 + 1}{x^2 - 1}$ is taken to be: 


 $\textbf{(A)}\ - 2 \qquad \textbf{(B)}\ 0 \qquad \textbf{(C)}\ \frac {3}{2} \qquad \textbf{(D)}\ \infty \qquad \textbf{(E)}\ -\frac{3}{2}$

 \textbf{Answer: }E



\textbf{Solution}

Factoring the numerator using the sum of cubes identity and denominator using the difference of squares identity gives 
 \[\dfrac{(x+1)(x^{2}-x+1)}{(x+1)(x-1)}\]Cancelling out a factor of $x+1$ from the numerator and denominator gives 
 \[\dfrac{(x^{2}-x+1)}{(x-1)}\]Plugging in $x= -1$ gives $\dfrac{3}{-2}$ or $\fbox{E}$.

 


	\item (\textit{AHSME 1952}) Two equal parallel chords are drawn $8$ inches apart in a circle of radius $8$ inches. The area of that part of the circle that lies between the chords is:


 $\textbf{(A)}\ 21\frac{1}{3}\pi-32\sqrt{3}\qquad \textbf{(B)}\ 32\sqrt{3}+21\frac{1}{3}\pi\qquad \textbf{(C)}\ 32\sqrt{3}+42\frac{2}{3}\pi \qquad\\ \textbf{(D)}\ 16\sqrt {3} + 42\frac {2}{3}\pi \qquad \textbf{(E)}\ 42\frac {2}{3}\pi$

 \textbf{Answer: }B



\textbf{Solution 1}


                    \begin{figure}[H]
                        \centering
                        \includegraphics[width=0.5\textwidth]{Solution/AHSME_I_1952_37_Solution_Figure_1.png}
                    \end{figure}
                    Draw in the diameter through A perpendicular to chords BD and CE. Label the intersection points of the diameter with the chords F and G. Then, AF=4, AG=4, and AFB ,AFD, AEG. and AGC are right triangles. We're going to find the area inside sectors BAD and CAE but outside triangles BAD and CAE.

 By pythagorean theorem, BF=$4\sqrt{3}$, as are DF, EG, and GC.

 It then follows that the area of triangles BAD and CAE are $16\sqrt{3}$. 

 Since AFB and AFD are 30-60-90 triangles, the area of sector BAD is $\frac{64\pi}{3}$, as is sector CAE.

 Thus, the area outside of the two chords is $\frac{128\pi}{3}-32\sqrt{3}$.Since we want the area inside the two chords, you can subtract the outside from the whole circle, which is $64\pi-(\frac{128\pi}{3}-32\sqrt{3})$ $\Rightarrow \frac{64\pi}{3}+32\sqrt{3}$, or $\fbox{B}$.

 


 



\textbf{Solution 2}

The area inside the chords are the arcs plus the triangles. We know that $\triangle AFB$ and $\triangle AFD$ are 30-60-90 triangles, which means that $\angle BAC = \angle DAE = 60^\circ$. This means that the sum of the sectors is $\frac{64 \pi} {3}$. The triangles are both equal to $\frac{8\sqrt{3} \cdot 4} {2}$, so the sum of the two is $32\sqrt{3}$. Finally, we add the two together. This yields the answer, $\frac{64\pi}{3}+32\sqrt{3}$, or $\fbox{B}$.

 


 



\textbf{Solution 3}

\textbf{Step 1: Finding the Height}

 We have a triangle with a slant side of length $8$ and a base half-length of $4$. By the Pythagorean theorem, the height $h$ satisfies $h^2 = 8^2 - 4^2$:

 
 \[h = \sqrt{8^2 - 4^2} = \sqrt{64 - 16} = \sqrt{48} = \sqrt{16 \cdot 3} = 4\sqrt{3}.\]

 \textbf{Step 2: Finding the Angles}

 Using the cosine ratio between the half-base $b = 4$ and the slant side $8$:
 \[\cos\theta = \frac{4}{8} = \frac{1}{2} \;\Longrightarrow\; \theta = 60^\circ.\]The full apex angle is $2\theta = 120^\circ$, so the sector angle is
 \[\theta' = 180^\circ - 120^\circ = 60^\circ.\]

 \textbf{Step 3: Area of the Triangles}

 Four congruent triangles, each base $4$ and height $4\sqrt{3}$:
 \[\text{Area}_\triangle = 4 \times \frac{1}{2} \times 4\sqrt{3} \times 4 = 32\sqrt{3}.\]

 \textbf{Step 4: Area of the Sectors}

 Each sector has $r = 8$ and spans $60^\circ = \frac{1}{6}$ of a circle:
 \[2 \times \frac{1}{6} \times \pi (8)^2 = \frac{128\pi}{6} = \frac{64\pi}{3}.\]

 \textbf{Step 5: Final Total Area}

 
 \[\boxed{32\sqrt{3} + \frac{64\pi}{3}}\]



	\item (\textit{AHSME 1952}) The area of a trapezoidal field is $1400$ square yards. Its altitude is $50$ yards. Find the two bases, if the number of yards in each base is an integer divisible by $8$. The number of solutions to this problem is: 


 $\textbf{(A)}\ \text{none} \qquad \textbf{(B)}\ \text{one} \qquad \textbf{(C)}\ \text{two} \qquad \textbf{(D)}\ \text{three} \qquad \textbf{(E)}\ \text{more than three}$

 \textbf{Answer: }D



\textbf{Solution}

Let us denote $8m$ and $8n$ to be our bases. Without loss of generality, let $m \le n$.

 Thus,
 \[50 * \frac{8m + 8n}{2} = 1400\]
 \[4m + 4n = 28\]
 \[m + n = 7\]

 Since $m$ & $n$ are integers, we see that the only solutions to this equation are $(1,6)$, $(2,5)$, and $(3,4)$. Therefore, the answer is $\fbox{(D) three}$



	\item (\textit{AHSME 1952}) If the perimeter of a rectangle is $p$ and its diagonal is $d$, the difference between the length and width of the rectangle is: 


 $\textbf{(A)}\ \frac {\sqrt {8d^2 - p^2}}{2} \qquad \textbf{(B)}\ \frac {\sqrt {8d^2 + p^2}}{2} \qquad \textbf{(C)}\ \frac{\sqrt{6d^2-p^2}}{2}\qquad\\ \textbf{(D)}\ \frac {\sqrt {6d^2 + p^2}}{2} \qquad \textbf{(E)}\ \frac {8d^2 - p^2}{4}$

 \textbf{Answer: }A



\textbf{Solution}


                    \begin{figure}[H]
                        \centering
                        \includegraphics[width=0.5\textwidth]{Solution/AHSME_I_1952_39_Solution_Figure_1.png}
                    \end{figure}
                    Let the sides of the rectangle be x and y. WLOG, assume x>y. Then, $2x+2y=p \Rightarrow x+y=\frac{p}{2}$.

 By pythagorean theorem, $x^2 + y^2 =d^2$.

 Since $x+y=\frac{p}{2}$, $(x+y)^2=\frac{p^2}{4} \Rightarrow x^2+2xy+y^2=\frac{p^2}{4}$.

 Rearranging to solve for $2xy$ gives $2xy = \frac{p^2}{4}-d^2$.

 Rearranging $(x-y)^2$ in terms of the defined variables becomes $(x-y)^2 = d^2 - (\frac{p^2}{4}-d^2)$. 

 In order to get (x-y), we have to take the square root of the expression and simplify. 

 $(x-y)=\sqrt{2d^2-\frac{p^2}{4}} \Rightarrow (x-y)=\sqrt{\frac{8d^2-p^2}{4}}$ $\Rightarrow$ $(x-y)=\frac{\sqrt{8d^2+p^2}}{2}$ $\Rightarrow$$\fbox{A}$.

 


	\item (\textit{AHSME 1952}) In order to draw a graph of $ax^2+bx+c$, a table of values was constructed. These values of the function for a set of equally spaced increasing values of $x$ were $3844, 3969, 4096, 4227, 4356, 4489, 4624$, and $4761$. The one which is incorrect is: 


 $\text{(A) } 4096 \qquad \text{(B) } 4356 \qquad \text{(C) } 4489 \qquad \text{(D) } 4761 \qquad \text{(E) } \text{none of these}$

 \textbf{Answer: }E



\textbf{Solution}

Since the polynomial is quadratic, its second differences must be constant. Taking the first differences, we have
 \begin{align*} 3969-3844&=125 \\ 4096-3969&=127 \\ 4227-4096&=131 \\ 4356-4227&=129 \\ 4489-4356&=133 \\ 4624-4489&=135 \\ 4761-4624&=137 \end{align*}
This leads to a common second difference of $2$, with the only discrepancy around the point $4227$. Observe that if this point were instead $4225$, the common second difference would, indeed be $2$ for all data points. Therefore the answer is $4227$, or $\fbox{E}$



	\item (\textit{AHSME 1952}) Increasing the radius of a cylinder by $6$ units increased the volume by $y$ cubic units. Increasing the altitude of the cylinder by $6$ units also increases the volume by $y$ cubic units. If the original altitude is $2$, then the original radius is: 


 $\text{(A) } 2 \qquad \text{(B) } 4 \qquad \text{(C) } 6 \qquad \text{(D) } 6\pi \qquad \text{(E) } 8$

 \textbf{Answer: }C



\textbf{Solution 1}

We know that the volume of a cylinder is equal to $\pi r^2h$, where $r$ and $h$ are the radius and height, respectively. So we know that $2\pi (r+6)^2-2\pi r^2=y=\pi r^2(2+6)-2\pi r^2$. Expanding and rearranging, we get that $2\pi (12r+36)=6\pi r^2$. Divide both sides by $6\pi$ to get that $4r+12=r^2$, and rearrange to see that $r^2-4r-12=0$. This factors to become $(r-6)(r+2)=0$, so $r=6$ or $r=-2$. Obviously, the radius cannot be negative, so our answer is $\fbox{(C) 6}$



	\item (\textit{AHSME 1952}) Let $D$ represent a repeating decimal. If $P$ denotes the $r$ figures of $D$ which do not repeat themselves, and $Q$ denotes the $s$ figures of $D$ which do repeat themselves, then the incorrect expression is:


 $\textbf{(A) } D = .PQQQ\ldots \qquad\\ \textbf{(B) } 10^rD = P.QQQ\ldots \\ \textbf{(C) } 10^{r + s}D = PQ.QQQ\ldots \qquad\\ \textbf{(D) } 10^r(10^s - 1)D = Q(P - 1) \\ \textbf{(E) } 10^r\cdot10^{2s}D = PQQ.QQQ\ldots$

 \textbf{Answer: }D



\textbf{Solution}

$\fbox{D}$



	\item (\textit{AHSME 1952}) The diameter of a circle is divided into $n$ equal parts. On each part a semicircle is constructed. As $n$ becomes very large, the sum of the lengths of the arcs of the semicircles approaches a length: 


 $\textbf{(A) } \qquad$ equal to the semi-circumference of the original circle


 $\textbf{(B) } \qquad$ equal to the diameter of the original circle


 $\textbf{(C) }$ greater than the diameter, but less than the semi-circumference of the original circle


 $\textbf{(D) } \qquad$ that is infinite


 $\textbf{(E) }$ greater than the semi-circumference


 
 \textbf{Answer: }A



\textbf{Solution}

Note that the half the circumference of a circle with diameter $d$ is $\frac{\pi*d}{2}$. 

 Let's call the diameter of the circle D. Dividing the circle's diameter into n parts means that each semicircle has diameter $\frac{D}{n}$, and thus each semicircle measures $\frac{D*pi}{n*2}$. The total sum of those is $n*\frac{D*pi}{n*2}=\frac{D*pi}{2}$, and since that is the exact expression for the semi-circumference of the original circle, the answer is $\boxed{A}$.

 


	\item (\textit{AHSME 1952}) If an integer of two digits is $k$ times the sum of its digits, the number formed by interchanging the the digits is the sum of the digits multiplied by


 $\textbf{(A) \ } 9-k  \qquad \textbf{(B) \ } 10-k \qquad \textbf{(C) \ } 11-k \qquad \textbf{(D) \ } k-1 \qquad \textbf{(E) \ } k+1$

 \textbf{Answer: }C



\textbf{Solution}

Let $n = 10a+b$. The problem states that $10a+b=k(a+b)$. We want to find $x$, where $10b+a=x(a+b)$. Adding these two equations gives $11(a+b) = (k+x)(a+b)$. Because $a+b \neq 0$, we have $11 = k + x$, or $x = \boxed{\textbf{(C) \ } 11-k}$.

 


	\item (\textit{AHSME 1952}) If $a$ and $b$ are two unequal positive numbers, then: 


 $\text{(A) } \frac{2ab}{a+b}>\sqrt{ab}>\frac{a+b}{2}\qquad \text{(B) } \sqrt{ab}>\frac{2ab}{a+b}>\frac{a+b}{2} \\ \text{(C) } \frac{2ab}{a+b}>\frac{a+b}{2}>\sqrt{ab}\qquad \text{(D) } \frac{a+b}{2}>\frac{2ab}{a+b}>\sqrt{ab} \\ \text{(E) } \frac {a + b}{2} > \sqrt {ab} > \frac {2ab}{a + b}$

 \textbf{Answer: }E



\textbf{Solution}

Using the RMS-AM-GM-HM inequality, we can see that the answer is $\fbox{E}$.

 


	\item (\textit{AHSME 1952}) The base of a new rectangle equals the sum of the diagonal and the greater side of a given rectangle, while the altitude of the new rectangle equals the difference of the diagonal and the greater side of the given rectangle. The area of the new rectangle is: 


 $\text{(A) greater than the area of the given rectangle} \quad\\ \text{(B) equal to the area of the given rectangle}  \quad\\ \text{(C) equal to the area of a square with its side equal to the smaller side of the given rectangle}  \quad\\ \text{(D) equal to the area of a square with its side equal to the greater side of the given rectangle}  \quad\\ \text{(E) equal to the area of a rectangle whose dimensions are the diagonal and the shorter side of the given rectangle}$

 \textbf{Answer: }C



\textbf{Solution}

Let the larger side of the rectangle be y and the smaller side of the rectangle be x. Then, by pythagorean theorem, the diagonal of the rectangle has length $\sqrt{x^{2}+y^{2}}$. 

 By the definition of the problem, the area of the new rectangle is $(\sqrt{x^{2}+y^{2}}+y)(\sqrt{x^{2}+y^{2}}-y)$

 Expanding this gives the area of the new rectangle to be $x^{2}$, or $\fbox{C}$.

 


	\item (\textit{AHSME 1952}) In the set of equations $z^x = y^{2x},\quad  2^z = 2\cdot4^x, \quad x + y + z = 16$, the integral roots in the order $x,y,z$ are: 


 $\textbf{(A) } 3,4,9 \qquad \textbf{(B) } 9,-5,-12 \qquad \textbf{(C) } 12,-5,9 \qquad \textbf{(D) } 4,3,9 \qquad \textbf{(E) } 4,9,3$

 \textbf{Answer: }D



\textbf{Solution #1}

The easiest method, which in this case is not very time consuming, is to guess and check and use process of elimination. We can immediately rule out (B) since it does not satisfy the third equation. The first equation allows us to eliminate (A) and (C) because the bases are different and aren't powers of each other.Finally, we manually check (D) and (E) and see that our answer is $\fbox{D}$




\textbf{Solution #2}

Let us actually solve the problem through algebra. We start by simplifying each of the equations given to us.

 Equation \#1
 \[z^x = y^{2x}\]
 \[z = y^2\]

 Equation \#2
 \[2^z = 2 * 4^x\]
 \[2^z = 2^{2x + 1}\]
 \[z = 2x + 1\]

 Equation \#3
 \[x + y + z = 16\]

 Equations 1 & 2 tell us that $z$ is an odd perfect square. They also tell us that $z$ is the largest of the three integers. We see that if we assume $z = 9$, then $y = 3$ and $x = 4$, which satisfies all three equations.

 Therefore, the answer is $\fbox{(D) 4,3,9}$



	\item (\textit{AHSME 1952}) Two cyclists, $k$ miles apart, and starting at the same time, would be together in $r$ hours if they traveled in the same direction, but would pass each other in $t$ hours if they traveled in opposite directions. The ratio of the speed of the faster cyclist to that of the slower is: 


 $\text{(A) } \frac {r + t}{r - t} \qquad \text{(B) } \frac {r}{r - t} \qquad \text{(C) } \frac {r + t}{r} \qquad \text{(D) } \frac{r}{t}\qquad \text{(E) } \frac{r+k}{t-k}$

 \textbf{Answer: }A



\textbf{Solution}


                    \begin{figure}[H]
                        \centering
                        \includegraphics[width=0.5\textwidth]{Solution/AHSME_I_1952_48_Solution_Figure_1.png}
                    \end{figure}
                    




\textbf{Solution 2}

We have a simple formula $d = rt$. The times the problem gives us, $r$ and $t$ are kind of annoying, so I will just let $x$ and $y$ be $r$ and $t$, respectively. I will also let $r_{1}$ and $r_{2}$ being the rate of the riders, where $r_{1}>r_{2}$. We can then write our equations based off these numbers, so I got $k = (r_1-r_2)x$ and $k = (r_1+r_2)y$. Rearranging the equations and solving for $r_1$ and $r_2$, I got $r_1 = k/x + r_2$ (from first equation) and $r_1 = k/y - r_2$ (from second equation). Also, we have $r_2 = k/y - r_1$ and $r_2 = -(k/x) + r_1$. Adding the first two together and the last two together, I got $2r_1 = \frac{(xk+yk)}{xy}$ and $2r_2 = \frac{(xk - yk)}{xy}$. Finally, dividing $2r_1$ by $2r_2$ gives us $r_1/r_2 = \frac{(x+y)}{(x-y)}$. Substituting $x$ and $y$ for $r$ and $t$, I got $\boxed{\text{(A) } \frac {r + t}{r - t}}$ -DragonFish12345 with credits to @RJ5303707 for LATEX

 


	\item (\textit{AHSME 1952}) 
                    \begin{figure}[H]
                        \centering
                        \includegraphics[width=0.5\textwidth]{Problem/AHSME_I_1952_49_Problem_Figure_1.png}
                    \end{figure}
                    


 In the figure, $\overline{CD}$, $\overline{AE}$ and $\overline{BF}$ are one-third of their respective sides. It follows that $\overline{AN_2}: \overline{N_2N_1}: \overline{N_1D} = 3: 3: 1$, and similarly for lines BE and CF. Then the area of triangle $N_1N_2N_3$ is: 


 $\text{(A) } \frac {1}{10} \triangle ABC \qquad \text{(B) } \frac {1}{9} \triangle ABC \qquad \text{(C) } \frac{1}{7}\triangle ABC\qquad \text{(D) } \frac{1}{6}\triangle ABC\qquad \text{(E) } \text{none of these}$

 \textbf{Answer: }C



\textbf{Solution}

Let $[ABC]=K.$ Then $[ADC] = \frac{1}{3}K,$ and hence $[N_1DC] = \frac{1}{7} [ADC] = \frac{1}{21}K.$ Similarly, $[N_2EA]=[N_3FB] = \frac{1}{21}K.$ Then $[N_2N_1CE] = [ADC] - [N_1DC]-[N_2EA] = \frac{5}{21}K,$ and same for the other quadrilaterals. Then $[N_1N_2N_3]$ is just $[ABC]$ minus all the other regions we just computed. That is, 
 \[[N_1N_2N_3] = K - 3\left(\frac{1}{21}K\right) - 3\left(\frac{5}{21}\right)K = K - \frac{6}{7}K = \boxed{\textbf{(C) }\frac{1}{7}[ABC]}\]

 \textbf{Alternative but very similar Solution} Let $[ABC]=K.$ Then $[ADC] = \frac{1}{3}K,$ and hence $[N_1DC] = \frac{1}{7} [ADC] = \frac{1}{21}K.$ Similarly, $[N_2EA]=[N_3FB]=[N_1DC] = \frac{1}{21}K.$ $[ADC]+[BEA]+[CFB]=\frac{1}{3}K+\frac{1}{3}K+\frac{1}{3}K = K.$ Then we can implement a similar but different area addition postulate to the first solution. It will be $[ADC]+[BEA]+[CFB]=K-\frac{1}{21}K-\frac{1}{21}K-\frac{1}{21}K+[N_1N_2N_3]$ (PIE in action). Using transitive property $K-\frac{1}{21}K-\frac{1}{21}K-\frac{1}{21}K+[N_1N_2N_3] = K.$ Subtracting and adding on both sides gives: 
 \[[N_1N_2N_3] = \boxed{\textbf{(C) }\frac{1}{7}[ABC]}\] ~many credits to the first solution ~Lopkiloinm

 



\textbf{Solution 2 (best solution)}

We can force this triangle to be equilateral because the ratios are always $3:3:1$ no matter which rotation, and with the symmetry of the equilateral triangle, it is safe to assume that this may be the truth (symmetry is very important, kids). Then, we can do a simple coordinate bash. Let $B$ be at $(0,0)$, $A$ be at $(3,3\sqrt{3})$, and $C$ be at $(6,0)$. We then create a new point $O$ at the center of everything. It should be noted because of similarity between $\triangle N_{1}N_{2}N_{3}$ and $\triangle ABC$, we can find the scale factor between the two triangle by simply dividing $N_{2}O$(not nitrous oxide) by $AO$. First, we need to find the coordinates of $O$ and $N_{2}$. $O$ is easily found at $(3,\sqrt{3})$ and $N_{2}$ be found by calculating equation of $BE$ and $AD$.$E$ is located $(4,2\sqrt{3})$ so $BE$ is $y=\frac{x\sqrt{3}}{2}$. $D$ be at $(4,0)$ and the slope is $-3\sqrt{3}$. We see that they be at the same $x$-value. Quick maths calculate the x value to be $4-\frac{2\sqrt{3}}{3\sqrt{3}+\frac{\sqrt{3}}{2}}$ which be $3\frac{3}{7}$. Another quick maths caculation of the $y$-value lead it be equal $2\sqrt{3}*\frac{3\sqrt{3}}{\frac{\sqrt{3}}{2}+3\sqrt{3}}$ which be $1\frac{5}{7}\sqrt{3}$. Peferct, so now $N_{2}$ be at $(3\frac{3}{7},1\frac{5}{7}\sqrt{3})$. Subtracting this coordinate with the coordinates of the center $O$ give you $(\frac{3}{7}, \frac{5}{7}\sqrt{3})$. Pythagorean theorem on that gives $\frac{\sqrt{84}}{7}$ which simplifies to $\frac{2\sqrt{3}}{\sqrt{7}}.$ $A$ is vertically above $O$ by $2\sqrt{3}$ units. The scale factor is thus $\frac{1}{\sqrt{7}}$ giving us an answer of: ~Lopkiloinm 
 \[\boxed{\textbf{(C) }\frac{1}{7}[ABC]}\] (Note: the presence of $7$ in the denominator gives hints on the answer, so when you see it, 1/7 looks like the obvious choice) 
 \[\](Another note: the question gives you the ratio of $3:3:1$ so I did not need to use that many steps to calculate the coordinate of $N_2$, directly attaining it to be at $(3\frac{3}{7},1\frac{5}{7}\sqrt{3})$ After realizing this, this method should actually be pretty quick, may even be quicker than the non coordinate method. Regardless, if you don't get the ratio and you encounter a problem similar to this, the coordinate bash with all these calculations should get you the correct answer. Also, there is an even bashier way using all the points and shoelace formula. That takes too much time.)

 



\textbf{Solution 3}

Visual graphical proof.\href{https://en.m.wikipedia.org/wiki/One-seventh\_area\_triangle\#/media/File\%3ATriangleOneSeventhAreaGraphicalSoln.png}{https://en.m.wikipedia.org/wiki/One-seventh\_area\_triangle\#/media/File\%3ATriangleOneSeventhAreaGraphicalSoln.png}

 



\textbf{Solution 4}

Using Routh's Theorem

 $x = \frac{BD}{CD} = \frac{\frac{2BC}{3}}{\frac{BC}{3}} = \frac{2}{1}$and$x = y = z$

 Hence

 $[N_1N_2N_3] = [ABC]\cdot\frac{(xyz-1)^2}{(xy+y+1)(yz+z+1)(zx+x+1)}$

 Substituting $x = y = z = 2$
$[N_1N_2N_3] = [ABC]\cdot\frac{7^2}{7^3}$$[N_1N_2N_3] = [ABC]\cdot\frac{1}{7}$

 Therefore the answer is,

 (C) $[N_1N_2N_3] = \frac{[ABC]}{7}$




\textbf{Solution 5 (Super fast)}

We can use triangle area ratios. $[N_1N_2N_3]=[ABD]-[BFN_3]-[N_3N_2AF]-[N_3N_1DB]$.

 So let $[ABC]=S$. Then $[ABD]=\frac{2}{3}S, [ADC]=\frac{1}{3}S, [N_1DC]=\frac{1}{21}S$.

 Similarly, $[FBC]=\frac{1}{3}S$. So $[FBDN_1]=\frac{6}{21}S$. Similarly, we can use this argument to find $[ABE]=\frac{1}{3}S, [FBN_3]=\frac{1}{21}S, [AN_2E]=\frac{1}{21}S$. So $[FN_3N_2A]=\frac{5}{21}S$. So $[N_1N_2N_3]=\frac{2}{3}S-\frac{5}{21}S-\frac{6}{21}S=\frac{1}{7}S$. 

 So therefore, since $[N_1N_2N_3]=\frac{1}{7}[ABC]$, select $\boxed{C}$.

 You can also use a similar idea (Principle of Inclusion-Exclusion) to find $[N_1N_2N_3]$. Notice that the triangle we desire is the intersection set $\triangle{AFC}, \triangle{ABD}, \triangle{BEC}$. So we just PIE it out.

 


	\item (\textit{AHSME 1952}) A line initially 1 inch long grows according to the following law, where the first term is the initial length. 
$1+\frac{1}{4}\sqrt{2}+\frac{1}{4}+\frac{1}{16}\sqrt{2}+\frac{1}{16}+\frac{1}{64}\sqrt{2}+\frac{1}{64}+\cdots$


 If the growth process continues forever, the limit of the length of the line is: 


 $\textbf{(A) } \infty\qquad \textbf{(B) } \frac{4}{3}\qquad \textbf{(C) } \frac{8}{3}\qquad \textbf{(D) } \frac{1}{3}(4+\sqrt{2})\qquad \textbf{(E) } \frac{2}{3}(4+\sqrt{2})$

 \textbf{Answer: }D



\textbf{Solution}

We can rewrite our sum as the sum of two infinite geometric sequences.
 \[1 + \frac{1}{4}\sqrt{2} + \frac{1}{4} + \frac{1}{16}\sqrt{2} + \frac{1}{16} + ... =\]
 \[(1 + \frac{1}{4} + \frac{1}{16} + \frac{1}{64} + ...) + (\frac{1}{4}\sqrt{2} + \frac{1}{16}\sqrt{2} + \frac{1}{64}\sqrt{2} + ...)\]
 \[(1 + \frac{1}{4} + \frac{1}{16} + \frac{1}{64} + ...) + \sqrt{2}(\frac{1}{4} + \frac{1}{16} + \frac{1}{64} + ...)\]We now take the sum of each of the infinite geometric sequences separately
 \[(1 + \frac{1}{4} + \frac{1}{16} + \frac{1}{64} + ...) + \sqrt{2}(\frac{1}{4} + \frac{1}{16} + \frac{1}{64} + ...) =\]
 \[\frac{1}{1 - \frac{1}{4}} + \sqrt{2}(\frac{\frac{1}{4}}{1 - \frac{1}{4}})\]
 \[\frac{4}{3} + \frac{\sqrt{2}}{3}\]
 \[\frac{1}{3}(4 + \sqrt{2})\]

 Therefore, the answer is $\fbox{(D)}$



\end{enumerate}